\chapter{Hinweise zur Verwendung von \LaTeX}
\label{cha:Hinweise_Latex}

Dieses Kapitel gibt allgemeine Hinweise zur Erstellung einer wissenschaftlichen Arbeit mit dem Textsatzprogramm \LaTeX.


\section{Vorlagen}

Die Arbeit soll auf dem Paket \verb|tuda-ci|, welches von der TU Darmstadt zur Verfügung gestellt wird, aufbauen.
Dies erleichtert es, sich an die Layout-Vorgaben sowie an die formalen Vorgaben (Selbständigkeitserklärung in der aktuellen Fassung) zu halten.

Basierend auf diesem Paket stellt das Fachgebiet rtm verschiedene Vorlagen zur Verfügung, um den Einstieg zu erleichtern:
\begin{itemize}
  \item \verb|tuda-ci_vorlage|\\
    Minimalbeispiel und Startpunkt für die eigene Arbeit.
    Bindet die notwendige Dokumentenklasse aus \verb|tuda-ci| ein, stellt das Fachgebietslogo zur Verfügung und gibt die minimale Grundstruktur der Arbeit vor.\\
    Im Quelltext ist die Anpassung auf verschiedene Varianten einer studentischen Arbeit (Proseminar, externe Arbeit in Unternehmen) erläutert.
  \item \verb|tuda-ci_thesisinfo|\\
    Quelltext zu diesem Dokument.
    Basiert ebenfalls auf \verb|tuda-ci| und zeigt exemplarisch, wie ein etwas größeres Dokument mit mehreren Dateien in einer Ordnerstruktur organisiert werden könnte.
    Man kann sich hieran orientieren, es ist aber nicht als Grundlage der eigenen Arbeit vorgesehen.
  \item \verb|tuda-ci_useful-packages|\\
    Hier sind (minimale) Konfigurationen und die Verwendung ein paar wahrscheinlich nützlicher Pakete (\zB \verb|tikz|) demonstriert.
\end{itemize}



\section{\TeX-System}

Ziel ist es, aus dem in \LaTeX{} geschriebenen "`Quelltext"' ein PDF zu erzeugen. Hierzu wird ein \TeX-Programm, quasi der "`Compiler"', benötigt.
\begin{itemize}
  \item pdflatex\\
    Das meist am Institut verwendete Programm.
    Erzeugt direkt eine \verb|pdf|-Datei und kann dabei direkt \verb|png|- und \verb|jpg|-Bilder verarbeiten.
  \item lualatex\\
    Etwas neueres Programm.
    Erzeugt auch direkt eine \verb|pdf|-Datei.
    Für große Dateien etwas langsamer.
  \item xelatex\\
    Weitere Alternative.
    Erzeugt auch direkt eine \verb|pdf|-Datei.
    (Kann direkt auf die Systemschriftarten zugreifen.)
  \item latex\\
    Das "`ursprüngliche"' Latex-Programm.
    Erzeugt eine Ausgabe im \verb|dvi|-Format, welches dann erst in eine \verb|ps|-Datei und diese dann über Ghostscript in eine \verb|pdf|-Datei umgewandelt werden muss.
    Sollte nicht mehr verwendet werden.
\end{itemize}
In der Anwendung sollten die ersten beiden Möglichkeiten kaum einen Unterschied machen.

Neben dem eigentlichen \LaTeX-Programm werden im Allgemeinen noch weitere Programme benötigt, \zB \verb|biber| zur Erzeugung des Literaturverzeichnisses.

Ein \LaTeX-Programm erzeugt bei einem Durchlauf (Aufruf) eine Menge an Hilfsdateien, die dann weitere Programme oder auch \LaTeX{} selber (in weiteren Durchläufen) verwenden, um beispielsweise Inhaltsverzeichnisse, Literaturverzeichnisse und Referenzierungen zu generieren.
\Dah typischerweise wäre eine Aufrufreihenfolge
\begin{verbatim}
  pdflatex -> biber -> pdflatex -> pdflatex
\end{verbatim}
Um dies zu vereinfachen, kann ein geeigneter Editor und/oder das Programm \verb|latexmk| verwendet werden.

Neben diesen auszuführenden Programmen existieren noch eine sehr große Menge an frei verfügbaren Paketen, die die verschiedensten Befehle zur Verfügung stellen.
(Erzeugung von Grafiken, besonderer Formelsatz, besondere Tabellen, aufrechte griechische Kleinbuchstaben, \ldots)

Zusammen bilden diese Bestandteile das \TeX-System.


\section{\LaTeX-Distribution}
\label{sec:Distribution}

Damit nicht alle oben genannten Komponenten selber und einzeln geladen und installiert werden müssen, sollte auf sogenannte \LaTeX-Distributionen zurückgegriffen werden.
Diese kümmern sich idealerweise um alles, und -- je nach Konfiguration -- laden diese auch noch fehlende Pakete automatisch herunter, wenn diese zum ersten Mal verwendet werden.

Welche Distribution verwendet wird, hängt vom persönlichen Geschmack ab.
Es besteht grundsätzlich die Wahl zwischen \Miktex{}, das vor allem unter Windows weit verbreitet ist und \texlive{}, das eher unter Linux und Mac Anwendung findet.
Daneben gibt es mittlerweile auch Cloud-Lösungen (overleaf), die sich besonders für Gruppenarbeiten eignen, da mehrere Personen gleichzeitig am gleichen Dokument arbeiten können.

Zuerst ist die Distribution herunterzuladen und zu installieren, wobei im Allgemeinen die Standardeinstellungen ausreichend sind.

Das sollte dann eigentlich schon ausreichen, um aus den zur Verfügung gestellten (Minimal-)Beispielen und einem geeigneten Editor (siehe unten) die \verb|pdf|-Dateien zu erzeugen.


\section{Editor}
Für \LaTeX{} gibt es eine Vielzahl von freien und kommerziellen Texteditoren.
Das \Texstudio ist ein freier open-source Editor, der sich bei uns am Institut bewährt hat.
Alternativ kann beispielsweise auch Visual Studio Code mit der "`LaTeX Workshop"'-Erweiterung verwendet werden.

Viele Editoren verfügen über eingebettete PDF-Reader \bzw arbeiten gut \zB mit SumatraPDF zusammen.
\Dah das erzeugte PDF kann im Reader geöffnet bleiben%
\footnote{Dies ist auch der Hauptgrund, weshalb von der Verwendung des Adobe-Acrobat-Readers abgeraten wird.}%
, und man kann aus dem Quellcode an die entsprechende Stelle im PDF und umgekehrt springen.




\section{Verwenden von \LaTeX-Paketen}

Durch das Einbinden von Zusatzpaketen kann \LaTeX\ angepasst und erweitert werden.
Die Pakete werden mit dem Befehl \verb|\usepackage{...}| eingebunden, \ggf können auch noch Optionen in \verb|[...]| angegeben werden.


\section{Definition eigener Befehle}

Die Möglichkeit eigene Befehle in \LaTeX{} zu definieren und zu verwenden erleichtert das Erstellen einer wissenschaftlichen Arbeit deutlich.
So dient ein eigener Befehl oft dazu, häufig verwendete Befehlsfolgen kürzer und schneller schreiben zu können.
Von zentraler Bedeutung ist außerdem, dass man diesen Befehl einfach ändern kann.
Hat man \zB alle Matrizen mit einem eigenen Befehl versehen, der diese fett formatiert, so lässt sich dies auch schnell für alle Matrizen wieder ändern.
Entscheidet man sich Matrizen mit einem Unterstrich zu kennzeichnen, so ist lediglich die Anpassung des entsprechenden Befehls notwendig.

Dies lässt sich auch auf Variablennamen übertragen.
Definiert man \zB für $\tilde{\hat{\mat{x}}}_\mathrm{b2}$ einen neuen Befehl \verb|\xb2|, so verkürzt sich zum einen der Schreibaufwand.
Zum anderen lässt sich selbst in der Endphase der Arbeit die Variable umbenennen, \zB in $\mat{z}_2$, indem lediglich der Befehl verändert wird.
